\chapter{Mới và Cũ}
\markboth{Mới và Cũ}{New and Old}

\section{Mới không phải lúc nào cũng hay hơn.}
\begin{truyen}{Mới không phải lúc nào cũng hay hơn.}{Newness is not a guarantee of value.}
\chuhoa{C}{hỉ vì thấy phương pháp mới đang được nhắc nhiều, anh vội bỏ hết cách dạy cũ từng có hiệu quả. Một thời gian sau, lớp học rối hơn mà chất lượng không khá hơn. Cái mới đáng quý khi nó tốt thật, không chỉ vì nó mới.}
\textit{(Because a new method was being widely praised, he quickly abandoned older teaching ways that had worked well. After some time, the classroom became more confused while quality did not improve. What is new is valuable when it is truly better, not merely because it is new.)}
\end{truyen}
\begin{baihoc}
Đừng thần tượng cái mới chỉ vì nó lạ mắt.
\textit{(Do not idolize the new merely because it looks fresh.)}
\end{baihoc}
\begin{ghinhoanh}
\textbf{Short English to remember:}

New is not always better.
\end{ghinhoanh}
\begin{tuhocnhanh}
1. Em có đang thích một điều chỉ vì nó mới không? \textit{(Are you liking something only because it is new?)}

2. Tiêu chí nào giúp em phân biệt mới thật sự tốt với mới chỉ để gây chú ý? \textit{(What criteria help you distinguish truly useful newness from attention-seeking novelty?)}
\end{tuhocnhanh}
\ngancach

\section{Cũ nhưng còn giá trị nên không cần vội bỏ.}
\begin{truyen}{Cũ nhưng còn giá trị nên không cần vội bỏ.}{Some old things keep a quiet strength.}
\chuhoa{B}{à vẫn giữ thói quen viết thư tay cho con cháu vào dịp đặc biệt. Giữa thời nhắn tin nhanh, những lá thư ấy chậm hơn nhưng ở lại lâu hơn. Không phải cái cũ nào cũng lỗi thời; có điều cũ là nơi chứa chiều sâu mà nhịp sống mới dễ làm rơi mất.}
\textit{(Grandmother still wrote letters by hand on special occasions. In an age of fast messages, those letters were slower yet stayed longer. Not everything old is outdated; some old things hold a depth that modern pace easily drops.)}
\end{truyen}
\begin{baihoc}
Cũ không đồng nghĩa với vô ích; nhiều cái cũ là nền cho điều mới đứng vững.
\textit{(Old does not mean useless; many old things are the foundation on which the new stands.)}
\end{baihoc}
\begin{ghinhoanh}
\textbf{Short English to remember:}

Some old things still carry depth.
\end{ghinhoanh}
\begin{tuhocnhanh}
1. Điều cũ nào trong gia đình hoặc lớp học của em vẫn còn đáng giữ? \textit{(Which old thing in your family or classroom is still worth keeping?)}

2. Em có đang bỏ một điều tốt chỉ vì nó không thời thượng không? \textit{(Are you abandoning something good simply because it is not trendy?)}
\end{tuhocnhanh}
\ngancach

\section{Muốn đổi mới nhưng chưa hiểu nền cũ.}
\begin{truyen}{Muốn đổi mới nhưng chưa hiểu nền cũ.}{Innovation without roots becomes noise.}
\chuhoa{M}{ột nhóm học sinh muốn làm dự án khác hẳn mọi năm, nhưng chưa chịu đọc kỹ những gì khóa trước đã làm được và vấp phải. Vì thiếu hiểu nền cũ, cái mới của các em phải trả giá bằng nhiều lỗi lặp lại. Muốn đi xa hơn, đôi khi phải cúi xuống học lại điều đã có trước mình.}
\textit{(A student group wanted to create a project completely different from previous years, yet refused to study what earlier groups had achieved and where they had failed. Because the old foundation was ignored, the new work paid the price by repeating old mistakes. To go farther, we sometimes must bow down and learn what came before us.)}
\end{truyen}
\begin{baihoc}
Cái mới bền nhất thường mọc lên từ sự hiểu biết sâu về cái cũ.
\textit{(The strongest newness usually grows from deep understanding of what came before.)}
\end{baihoc}
\begin{ghinhoanh}
\textbf{Short English to remember:}

Good innovation knows its roots.
\end{ghinhoanh}
\begin{tuhocnhanh}
1. Em có đang muốn đổi quá nhanh mà chưa hiểu kỹ nền tảng không? \textit{(Are you trying to change too quickly without understanding the foundation?)}

2. Điều cũ nào em cần học lại để làm cái mới tốt hơn? \textit{(What older lesson do you need to revisit in order to do the new thing better?)}
\end{tuhocnhanh}
\ngancach

\section{Giữ một thói quen cũ tốt giữa nhiều cám dỗ mới.}
\begin{truyen}{Giữ một thói quen cũ tốt giữa nhiều cám dỗ mới.}{Old discipline can protect a new world.}
\chuhoa{E}{m vẫn giữ thói quen đọc sách giấy mỗi tối dù điện thoại có quá nhiều điều hấp dẫn. Chính thói quen tưởng chừng cũ ấy giúp em không bị kéo trôi hoàn toàn theo nhịp cuốn của màn hình. Có những cái cũ không làm mình chậm lại; nó giữ mình khỏi bị phân mảnh.}
\textit{(The student still kept the habit of reading printed books every night despite the many attractions of the phone. That seemingly old habit prevented total drift into the pull of screens. Some old things do not slow us down; they keep us from falling apart.)}
\end{truyen}
\begin{baihoc}
Một thói quen cũ tốt có thể là hàng rào quý trong thời đại quá nhiều kích thích mới.
\textit{(A good old habit can be a precious fence in an age of too many new stimulations.)}
\end{baihoc}
\begin{ghinhoanh}
\textbf{Short English to remember:}

Old discipline can guard your mind.
\end{ghinhoanh}
\begin{tuhocnhanh}
1. Thói quen cũ tốt nào em đang bỏ dở? \textit{(Which good old habit are you neglecting?)}

2. Em cần giữ điều gì để không bị cuốn trôi bởi quá nhiều cái mới? \textit{(What do you need to keep in order not to be swept away by too much novelty?)}
\end{tuhocnhanh}
\ngancach

\section{Làm mới điều cũ để nó tiếp tục sống.}
\begin{truyen}{Làm mới điều cũ để nó tiếp tục sống.}{Renewal can honor tradition.}
\chuhoa{K}{hông phải giữ cái cũ là giữ nguyên mọi hình thức của nó. Một cô giáo vẫn dạy những giá trị quen thuộc như biết ơn, trung thực, tự học, nhưng kể bằng ví dụ gần với đời sống học sinh hôm nay. Điều cũ không chết khi được làm mới đúng cách; nó chỉ đổi áo để tiếp tục đi cùng thời đại.}
\textit{(Keeping something old does not mean preserving every outer form unchanged. One teacher still taught familiar values such as gratitude, honesty, and self-learning, but through examples close to today's students. The old does not die when renewed well; it simply changes clothes to keep walking with the times.)}
\end{truyen}
\begin{baihoc}
Trung thành với giá trị không có nghĩa là bất động trong cách thể hiện.
\textit{(Being faithful to a value does not mean remaining motionless in the way it is expressed.)}
\end{baihoc}
\begin{ghinhoanh}
\textbf{Short English to remember:}

Renew the form, keep the value.
\end{ghinhoanh}
\begin{tuhocnhanh}
1. Điều cũ nào đáng được làm mới thay vì bỏ đi? \textit{(Which old thing deserves renewal instead of removal?)}

2. Em có giá trị nào muốn giữ nhưng cần cách thể hiện mới hơn? \textit{(Is there a value you want to keep but express in a fresher way?)}
\end{tuhocnhanh}
\ngancach
